\begin{abstract}
We present a self-consistent, spatially resolved cluster-dynamics model for the
irradiated microstructure of $\alpha$-zirconium, together with the construction
by which its continuum fields are converted into a discrete dislocation
population and the two-time-scale method by which the whole is integrated. The
model is self-consistent in two distinct senses. First, the bias factors of
\emph{all} variants follow from anisotropic diffusion theory: every capture
efficiency appearing in a loop growth law, and every sink strength those loops
impose on the mobile species, is generated from the \emph{same} diffusion
tensors that transport those species, through a single anisotropy parameter per
species, rather than from an independent set of phenomenological bias factors.
Second, the transfer between the continuum fields and the discrete loop
population preserves what defines that population --- stored defects exactly,
loop number to the integer draw, the dispersion of the size distribution to the
truncation of the sampling --- so a calculation may be run as continuum, as
discrete, or as a hybrid, with the choice made on grounds of dose and
computational cost rather than of physics.
The immobile population is resolved into nine crystallographic families --- the
three prismatic $\langle a\rangle$ variants of both interstitial and vacancy
character, and a three-state basal vacancy chain --- and each carries the first
three moments of its size distribution, closed by a log-normal ansatz, so that
the width of the distribution is transported and not only its mean. The mobile
species obey a steady reaction--diffusion problem with Dirichlet conditions at
grain boundaries; the immobile families obey pointwise ordinary differential
equations in which no neighbor's state appears. A Lie--Trotter operator split
exploits that structure: freezing the mobile field removes the spatial coupling
by construction, so that $7.8\times10^{7}$ stiff equations are integrated as
$72\,928$ independent $36$-dimensional systems rather than one coupled system of
$2.8\times10^{6}$, a reduction of $6.5\times10^{9}$ in the cost of a single
factorization. Grain boundaries are treated as sinks for both populations, by
different mechanisms: a Dirichlet condition for the mobile species, and for the
loops a size-gated removal of the fraction of each family's own size
distribution large enough to intersect the boundary. Applied to a
\SI{500}{\nano\meter} hexagonal single crystal at \SI{573}{\kelvin} to
\SI{10}{\dpaunit}, the model resolves mobile boundary layers of
$8$--$46$\,\si{\nano\meter} that order themselves by species mobility and narrow
by a factor of two as the sink density grows; closes both atom balances on
production, with the boundary absorbing $41\%$ of interstitial production as
mobile species and a further $5.1\%$ of vacancy production as loops; and
predicts a loop-denuded zone $\sim5$\,\si{\nano\meter} wide across which the mean
loop diameter recovers from $3.5$ to $58$\,\si{\nano\meter}. Carrying the loop
absorption channel removes an artifact that otherwise makes a domain average
meaningless: without it, the domain and interior means of the basal loop density
differ by a factor of $873$. The parameters are calibrated against 78
transmission-electron-microscope measurements of loop density and diameter as a
maximum \emph{a posteriori} estimate under stated priors, and we report
explicitly which of them the data determine and which rest on the bounds of the
admissible range. The model produces basal loops that are too numerous by a
factor of $9.5$ and too small by a factor of $1.9$; the two errors are linked by
the number and growth balances, and we argue that they identify a missing
channel that removes vacancy loop \emph{number} without shortening the
survivors' lives.
\end{abstract}

\begin{keyword}
zirconium \sep cluster dynamics \sep spatially resolved \sep dislocation loops
\sep diffusional anisotropy \sep grain boundary \sep irradiation damage
\end{keyword}
